\chapter{Modèles de volatilité}\label{ch:qm:volatility-models}

Un rapport de risque mesure la volatilité de chaque devise par l'écart-type à pondération égale de ses 250 derniers rendements quotidiens. Le 11 novembre
2022, l'euro a progressé de 3.5~\% face au dollar~: c'était le premier taux de référence de la BCE (fixé vers 14 h 10 CET) après l'indice des prix à la consommation américain d'octobre, publié
à 8 h 30, heure de New York, le 10 novembre, montrait un ralentissement de l'inflation annuelle à 7.7~\%. À la clôture du 2 novembre 2023, 250 fixages plus tard, ce rendement sort de la fenêtre, et la
volatilité EUR/USD que le rapport reporte au lendemain passe de 0.538~\% à 0.496~\%, soit une baisse de 7.8~\%, un jour où l'euro a progressé de 1.2~\%. Un modèle
qui laisserait s'estomper progressivement les anciennes nouvelles aurait affiché une hausse. Ce chapitre traite des modèles qui le font~: la famille ARCH et GARCH et leur
cas particulier à pondération exponentielle, la \omterm{def:qm:volatility-models:rv}{variance réalisée} à partir de données infrajournalières et l'autorégression hétérogène qui la prévoit, et les
fonctions de perte et tests qui tranchent quelle prévision est la meilleure lorsque la vraie variance n'est jamais observée.

\section{La famille des modèles à hétéroscédasticité conditionnelle}

\begin{definition}[Regroupement de volatilité, hétéroscédasticité conditionnelle]\label{def:qm:volatility-models:clustering}
Le \emph{regroupement de volatilité}\index{regroupement de volatilité} est la tendance des grands rendements, de signe quelconque, à être suivis de grands rendements, et les petits
suivis de petits~: les rendements au carré sont positivement autocorrélés, alors que les rendements eux-mêmes ne le sont presque pas. Une série de rendements présente une \emph{hétéroscédasticité
conditionnelle}\index{hétéroscédasticité conditionnelle} lorsque sa variance conditionnellement au passé, $h_t = \Var(r_t \mid \mathcal F_{t-1})$, varie dans le temps.
\end{definition}

\begin{definition}[Modèles ARCH et GARCH, persistance de la volatilité]\label{def:qm:volatility-models:garch}
Dans un \emph{modèle ARCH}\index{modèle ARCH} (Engle, 1982), $r_t = \sqrt{h_t}z_t$ avec des $z_t$ i.i.d. standardisés et $h_t = \omega + \sum_{i=1}^q\alpha_ir_{t-i}^2$.
Le \emph{modèle GARCH}\index{modèle GARCH} (Bollerslev, 1986) ajoute des variances retardées~; le GARCH(1,1) est $h_t = \omega + \alpha r_{t-1}^2 + \beta h_{t-1}$ avec $\omega > 0$,
$\alpha, \beta \ge 0$. Sa \emph{persistance de la volatilité}\index{persistance de la volatilité} est $\alpha + \beta$, soit le taux auquel un choc de variance s'amortit.
\end{definition}

\begin{proposition}[Stationnarité et demi-vie du GARCH(1,1)]\label{prop:qm:volatility-models:garch}
Si $\alpha + \beta < 1$, le rendement suivant un GARCH(1,1) est faiblement stationnaire, de variance inconditionnelle $\bar h = \omega/(1 - \alpha - \beta)$, et la
prévision de la variance à $k$ jours est $\E_t[h_{t+k}] = \bar h + (\alpha + \beta)^{k-1}(h_{t+1} - \bar h)$. Un choc sur la variance est divisé par deux en $\ln\frac12/\ln(\alpha + \beta)$
jours.
\end{proposition}

\begin{proof}
$\E_t[r_{t+k}^2] = \E_t[h_{t+k}]$, donc $\E_t[h_{t+k+1}] = \omega + (\alpha + \beta)\E_t[h_{t+k}]$, une récurrence linéaire de point fixe $\bar h$ et de taux $\alpha + \beta$.
En prenant les espérances inconditionnelles, on obtient $\E[r^2] = \bar h$ lorsque $\alpha + \beta < 1$.
\end{proof}

Les paramètres sont estimés en maximisant la vraisemblance des rendements conditionnellement à la récurrence, avec des innovations normales ou de Student-$t$. Lorsque la
vraisemblance normale est utilisée alors que les innovations ne sont pas normales, l'\omterm{def:qm:estimation:estimator}{estimateur} est dit de \omterm{def:qm:estimation:kl}{quasi-maximum de vraisemblance}~: il reste convergent et
asymptotiquement normal sous des conditions de régularité, et ses \omterm{def:qm:estimation:estimator}{erreurs types} doivent être calculées par le sandwich du chapitre~11.

\begin{definition}[Volatilité EWMA]\label{def:qm:volatility-models:ewma}
La \emph{volatilité EWMA}\index{volatilité EWMA} de facteur de décroissance $\lambda$ est $h_t = \lambda h_{t-1} + (1 - \lambda)r_{t-1}^2$~: une moyenne exponentiellement pondérée des carrés
des rendements passés. RiskMetrics (1996) a fixé $\lambda = 0.94$ pour les données quotidiennes.
\end{definition}

L'EWMA est un GARCH(1,1) avec $\omega = 0$, $\alpha = 1 - \lambda$ et $\beta = \lambda$~: intégré, $\alpha + \beta = 1$, il n'a donc pas de variance inconditionnelle vers laquelle revenir, et
sa prévision à tout horizon est celle d'aujourd'hui. Avec $\lambda = 0.94$, le poids d'un rendement est divisé par deux en 11 jours. La fenêtre à pondération égale est l'autre extrême~:
chaque rendement de la fenêtre a le même poids, et au-delà, aucun.

\begin{definition}[Effet de levier, modèle GJR-GARCH]\label{def:qm:volatility-models:gjr}
L'\emph{effet de levier}\index{effet de levier} est la tendance de la volatilité à augmenter davantage après des baisses qu'après des hausses de même ampleur, particulièrement marqué pour les
indices actions. Le \emph{modèle GJR-GARCH}\index{modèle GJR-GARCH} (Glosten, Jagannathan et Runkle, 1993) le capture par $h_t = \omega + (\alpha + \gamma\mathbf 1_{r_{t-1} <
0})r_{t-1}^2 + \beta h_{t-1}$~; sa persistance est $\alpha + \beta + \gamma/2$ pour des innovations symétriques.
\end{definition}

Estimés sur 7\,098 rendements quotidiens EUR/USD (1999--2026, en pourcentage), les trois modèles donnent~:
\begin{center}
\footnotesize
\setlength{\tabcolsep}{4pt}
\begin{tabular}{@{}lccccccc@{}}
\toprule
 & $\omega$ & $\alpha$ & $\beta$ & $\gamma$ & $\nu$ & $\alpha + \beta + \gamma/2$ & \omterm{def:qm:stochastic-differential-equations:ou}{demi-vie} \\
\midrule
GARCH, gaussien & 0.0010 & 0.029 (0.003) & 0.968 (0.004) & --- & --- & 0.9973 & 255 \\
GARCH, Student $t$ & 0.0007 & 0.030 (0.003) & 0.968 (0.003) & --- & 7.2 (0.6) & 0.9984 & 445 \\
GJR, Student $t$ & 0.0006 & 0.027 (0.005) & 0.969 (0.004) & 0.006 (0.005) & 7.2 (0.6) & 0.9986 & 482 \\
\bottomrule
\end{tabular}
\end{center}
avec des \omterm{def:qm:estimation:estimator}{erreurs types} de sandwich entre parenthèses. La vraisemblance de Student-$t$ est supérieure de 141 points de log-vraisemblance pour un paramètre supplémentaire~; l'asymétrie n'est pas
significative (une paire de devises n'a pas de direction naturelle «~baissière~», contrairement à un indice actions)~; et la persistance est si proche de 1 qu'un choc sur la volatilité EUR/USD
met bien plus d'un an à se réduire de moitié~: la \omterm{def:qm:linear-time-series:longmem}{mémoire longue} de la volatilité (chapitre~17), vue à travers un modèle à mémoire courte
(la \cref{fig:qm:volatility-models:history}).

\begin{omfigure}
\begin{tikzpicture}
\begin{axis}[width=12cm,height=5cm,xlabel={année},ylabel={\% quotidien (rendement absolu, volatilité)},
  tick label style={font=\small},label style={font=\small},xmin=1999,xmax=2026.8,ymin=0,ymax=3,grid=major,grid style={gray!25},
  xtick={2000,2005,2010,2015,2020,2025},xticklabel style={/pgf/number format/set thousands separator={}},
  legend style={font=\footnotesize,at={(0.98,0.96)},anchor=north east,fill=white,draw=none},table/col sep=comma]
\addplot[only marks,mark=*,mark size=0.3pt,black!35] table[x=year,y=absr]{figdata/methods/18-volatility-models/history.csv};
\addplot[omDef,very thick] table[x=year,y=garch]{figdata/methods/18-volatility-models/history.csv};
\legend{rendement quotidien absolu (un jour sur cinq),volatilité GARCH-$t$}
\end{axis}
\end{tikzpicture}

\omcaption{Rendements absolus quotidiens EUR/USD et volatilité conditionnelle du GARCH(1,1) estimé avec des innovations de Student-$t$, 1999--2026. Les
\omterm{def:qm:volatility-models:clustering}{regroupements de volatilité} (2000, 2008--2009, 2011, 2015, 2022) sont ce que décrit la persistance de 0.9984 du modèle. Données~: taux de référence de l'euro de la BCE (source~: statistiques
de la BCE).}\label{fig:qm:volatility-models:history}
\end{omfigure}

Les trois \omterm{def:qm:estimation:estimator}{estimateurs} traitent très différemment un grand rendement (la \cref{fig:qm:volatility-models:impulse}). Une journée à 3.5~\% ajoute à la prévision de variance un montant
qui décroît au rythme de 0.9984 par jour sous le GARCH estimé, de 0.94 par jour sous l'EWMA, et pas du tout sous la fenêtre, qui le conserve intact pendant 250
jours puis le supprime d'un coup. Ce retrait est le fantôme du rendement de l'accroche (la \cref{fig:qm:volatility-models:ghost})~: le jour où il s'est produit, la volatilité de la fenêtre a baissé
de 7.8~\%, tandis que la volatilité GARCH passait de 0.424~\% à 0.465~\% et l'EWMA de 0.410~\% à 0.490~\%, les deux réagissant au mouvement de 1.2~\% de la journée.

\begin{omfigure}
\begin{tikzpicture}
\begin{semilogyaxis}[width=10cm,height=5cm,xlabel={jours après un rendement de 3.5~\%},ylabel={variance supplémentaire ($\%^2$)},
  tick label style={font=\small},label style={font=\small},xmin=0,xmax=300,ymin=1e-4,ymax=1,grid=major,grid style={gray!25},
  legend columns=3,legend style={font=\footnotesize,at={(0.5,-0.3)},anchor=north,draw=none,fill=none,
  /tikz/every even column/.append style={column sep=8pt}},table/col sep=comma]
\addplot[omDef,very thick] table[x=day,y=garch]{figdata/methods/18-volatility-models/impulse.csv};
\addplot[omThm,very thick,dashed] table[x=day,y=ewma]{figdata/methods/18-volatility-models/impulse.csv};
\addplot[black,very thick,dotted,const plot] table[x=day,y=window]{figdata/methods/18-volatility-models/impulse.csv};
\legend{{GARCH-$t$, estimé},{EWMA, 0.94},fenêtre de 250 jours}
\end{semilogyaxis}
\end{tikzpicture}

\omcaption{Durée pendant laquelle un rendement quotidien de 3.5~\% reste présent dans chaque prévision de variance (la variance supplémentaire par rapport à un fond de journées à 0.5~\%). Le GARCH le laisse décroître lentement
(\omterm{def:qm:stochastic-differential-equations:ou}{demi-vie} de 445 jours), l'EWMA rapidement (11 jours), et la fenêtre à pondération égale le conserve intact pendant 250 jours puis le supprime en une seule fois. Données~: le
tutoriel du chapitre, à partir des paramètres estimés.}\label{fig:qm:volatility-models:impulse}
\end{omfigure}

\begin{omfigure}
\begin{tikzpicture}
 \begin{axis}[width=11cm,height=5cm,xlabel={jours de négociation par rapport au 2 novembre 2023},ylabel={volatilité quotidienne (\%)},
  tick label style={font=\small},label style={font=\small},xmin=-300,xmax=60,ymin=0.3,ymax=1.2,grid=major,grid style={gray!25},
  legend columns=3,legend style={font=\footnotesize,at={(0.5,-0.3)},anchor=north,draw=none,fill=none,
  /tikz/every even column/.append style={column sep=8pt}},table/col sep=comma]
\addplot[black,very thick,const plot] table[x=k,y=window]{figdata/methods/18-volatility-models/ghost.csv};
\addplot[omDef,very thick] table[x=k,y=garch]{figdata/methods/18-volatility-models/ghost.csv};
\addplot[omThm,thick,dashed] table[x=k,y=ewma]{figdata/methods/18-volatility-models/ghost.csv};
\legend{fenêtre de 250 jours,GARCH-$t$,{EWMA, 0.94}}
\end{axis}
\end{tikzpicture}

\omcaption{Prévisions de volatilité EUR/USD autour du jour où le rendement du 11 novembre 2022 est sorti de la fenêtre de 250 jours (jour 0, 2 novembre 2023). La fenêtre bondit
d'une même valeur lorsque le rendement entre (jour $-250$) et chute de 7.8~\% lorsqu'il sort~; les prévisions GARCH et EWMA réagissent au marché, et non au
calendrier. Données~: taux de change de référence de l'euro de la BCE (source~: statistiques de la BCE).}\label{fig:qm:volatility-models:ghost}
\end{omfigure}

\section{Volatilité réalisée}

\begin{definition}[Variance intégrée, variance réalisée]\label{def:qm:volatility-models:rv}
Si le logarithme du prix suit $dX_t = \mu_t\,dt + \sigma_t\,dW_t$, la \emph{variance intégrée}\index{variance intégrée} du jour $t$ est $\mathrm{IV}_t = \int_{t-1}^t\sigma_s^2\,ds$.
La \emph{variance réalisée}\index{variance réalisée} est la somme des carrés des rendements intrajournaliers du jour, $\mathrm{RV}_t = \sum_{j=1}^m r_{t,j}^2$ sur $m$ intervalles.
\end{definition}

\begin{proposition}[La variance réalisée converge vers la variance intégrée]\label{prop:qm:volatility-models:rv}
En l'absence de bruit de microstructure, $\mathrm{RV}_t \to \mathrm{IV}_t$ en probabilité lorsque $m \to \infty$~; si $\sigma$ est constant au cours de la journée, $\mathrm{RV}_t/\mathrm{IV}_t$ a
une moyenne égale à un et un écart type $\sqrt{2/m}$.
\end{proposition}

\begin{proof}
$\mathrm{RV}_t$ est la \omterm{def:qm:brownian-motion:qv}{variation quadratique} discrétisée de $X$, qui converge vers $\int\sigma^2$ (chapitre~3). Si $\sigma$ est constant, chaque $r_{t,j} \sim \mathcal N(0, \mathrm{IV}_t/m)$
(la dérive est négligeable à cette échelle), de sorte que $m\,\mathrm{RV}_t/\mathrm{IV}_t$ suit une loi $\chi^2_m$, de moyenne $m$ et de variance $2m$.
\end{proof}

Un rendement quotidien au carré est une \omterm{def:qm:volatility-models:rv}{variance réalisée} avec $m = 1$~: une mesure sans biais mais très bruitée de la variance du jour, avec une erreur relative de $\sqrt2 = 141\%$.
Les rendements à cinq minutes sur une séance de 6.5~heures ($m = 78$) la ramènent à 16~\%. Le marché simulé du chapitre comporte trois composantes de volatilité (coefficients AR quotidiens
0.3, 0.9 et 0.99)~; les erreurs relatives mesurées valent 144~\%, 16~\% et 7.3~\% pour $m = 1$, 78 et 390, contre 141~\%, 16~\% et 7.2~\% pour la théorie
(la \cref{fig:qm:volatility-models:rv}). Sur les marchés réels, le gain cesse à la fréquence où le rebond bid-ask et la discrétisation dominent les rendements, le
thème du chapitre~21.

\begin{omfigure}
\begin{tikzpicture}
\begin{loglogaxis}[width=9cm,height=5cm,xlabel={rendements intrajournaliers par jour, $m$},ylabel={écart type de RV/IV $- 1$},
  tick label style={font=\small},label style={font=\small},xmin=0.8,xmax=500,ymin=0.05,ymax=2,grid=major,grid style={gray!25},
  xtick={1,6,13,26,78,390},xticklabels={1,6,13,26,78,390},
  legend style={font=\footnotesize,at={(0.97,0.97)},anchor=north east,fill=white,draw=none},table/col sep=comma]
\addplot[only marks,mark=*,mark size=2pt,omDef] table[x=m,y=error]{figdata/methods/18-volatility-models/rv.csv};
\addplot[omThm,thick,dashed] table[x=m,y=theory]{figdata/methods/18-volatility-models/rv.csv};
\legend{marché simulé,$\sqrt{2/m}$}
\end{loglogaxis}
\end{tikzpicture}

\omcaption{Erreur relative de la \omterm{def:qm:volatility-models:rv}{variance réalisée} comme mesure de la \omterm{def:qm:volatility-models:rv}{variance intégrée}, en fonction du nombre de rendements intrajournaliers (1 = rendement quotidien au carré, 78 =
cinq minutes sur 6.5~heures, 390 = une minute), dans un marché simulé à volatilité stochastique à trois composantes, sans bruit de microstructure, et la théorie
$\sqrt{2/m}$. Données~: le tutoriel du chapitre, à graine fixée.}\label{fig:qm:volatility-models:rv}
\end{omfigure}

\section{Modèles autorégressifs hétérogènes}

\begin{definition}[Modèle HAR]\label{def:qm:volatility-models:har}
Le \emph{modèle HAR}\index{modèle HAR} (Corsi, 2009) régresse la \omterm{def:qm:volatility-models:rv}{variance réalisée} de demain sur celle d'aujourd'hui, sur la moyenne des cinq derniers jours et sur la moyenne des
vingt-deux derniers jours~: $\mathrm{RV}_{t+1} = b_0 + b_d\mathrm{RV}_t + b_w\mathrm{RV}_t^{(5)} + b_m\mathrm{RV}_t^{(22)} + e_{t+1}$.
\end{definition}

Trois horizons de moyennage, estimés par moindres carrés, reproduisent la décroissance lente et d'allure hyperbolique des autocorrélations de la volatilité avec un modèle AR d'ordre 22
restreint à trois paramètres~: une \omterm{def:qm:linear-time-series:longmem}{mémoire longue} approchée par une cascade de composantes courtes. Sur 2\,000 jours du marché simulé, l'ajustement donne $b_d = 0.35$, $b_w =
0.27$, $b_m = 0.20$ avec $R^2 = 0.37$~; sur les 1\,000 jours suivants, ses prévisions battent un AR(1) sur la \omterm{def:qm:volatility-models:rv}{variance réalisée} (QLIKE $-0.479$ contre $-0.473$,
statistique de Diebold--Mariano $-4.5$) et sont plus proches de la \omterm{def:qm:volatility-models:rv}{variance intégrée} réelle (erreur quadratique moyenne de 0.0073 contre 0.0079).

\section{Prévision et évaluation}

La variance d'un jour n'est jamais observée, seulement un proxy~: un rendement au carré ou une \omterm{def:qm:volatility-models:rv}{variance réalisée}. Une fonction de perte compare une prévision $h$ au proxy
$\hat\sigma^2$~; Patton (2011) a montré que le classement des prévisions n'est préservé, face à un proxy bruité mais sans biais, que pour une famille de pertes qui comprend
l'erreur quadratique et la \omterm{def:qm:volatility-models:eval}{perte QLIKE}.

\begin{definition}[Perte QLIKE, régression de Mincer--Zarnowitz, test de Diebold--Mariano]\label{def:qm:volatility-models:eval}
La \emph{perte QLIKE}\index{perte QLIKE} vaut $\hat\sigma^2/h + \ln h$ (à des termes près ne dépendant pas de $h$), l'opposé de la log-vraisemblance gaussienne du
rendement du proxy~: elle est minimisée en espérance par $h = \E[\hat\sigma^2]$ et pénalise davantage la sous-estimation que la surestimation. La \emph{régression de Mincer--Zarnowitz
}\index{régression de Mincer--Zarnowitz} (1969) du proxy sur la prévision, $\hat\sigma^2_t = a + bh_t + e_t$, teste l'absence de biais ($a = 0$, $b = 1$). Le
\emph{test de Diebold--Mariano}\index{test de Diebold--Mariano} (1995) de précision égale divise la différence des pertes moyennes de deux prévisions par son \omterm{def:qm:estimation:estimator}{erreur type} HAC.
\end{definition}

Sur EUR/USD, le GARCH-$t$ n'est réestimé que sur 1999--2014 (persistance 0.9982, \omterm{def:qm:stochastic-differential-equations:ou}{demi-vie} de 392 jours) et prolongé à paramètres fixes~; l'EWMA et la fenêtre
n'ont besoin d'aucun ajustement. Sur les 3\,002 jours de test de 2015 à 2026, avec le rendement au carré comme proxy, la QLIKE moyenne vaut $-0.521$ pour le GARCH, $-0.495$ pour l'EWMA
et $-0.435$ pour la fenêtre (la \cref{fig:qm:volatility-models:eval}). Le GARCH devance l'EWMA avec une statistique de Diebold--Mariano de $-2.76$ ($p = 0.006$) et la fenêtre
avec $-4.10$~; l'EWMA devance la fenêtre avec $-2.48$. Les pentes de Mincer--Zarnowitz valent 0.77, 0.60 et 0.56, toutes inférieures à un (chaque prévision surréagit aux
données récentes), et les $R^2$ valent 0.030, 0.027 et 0.010~: un rendement quotidien au carré est si bruité que même une bonne prévision n'en explique que 3~\%.

\begin{omfigure}
\begin{tikzpicture}
\begin{axis}[width=9cm,height=4.2cm,xbar,bar width=10pt,
  xlabel={QLIKE moyenne moins celle de la fenêtre de 250 jours (2015--2026)},
  ytick={0,1,2},yticklabels={GARCH-$t$,{EWMA, 0.94},fenêtre de 250 jours},
  tick label style={font=\small},label style={font=\small},
  xmin=-0.1,xmax=0.02,ymin=-0.6,ymax=2.6,grid=major,grid style={gray!25},
  xticklabel style={/pgf/number format/fixed,/pgf/number format/precision=2},scaled x ticks=false,table/col sep=comma]
\addplot[fill=omDef!60,draw=omDef,nodes near coords,nodes near coords style={font=\footnotesize,anchor=west,xshift=2pt,
  /pgf/number format/fixed,/pgf/number format/precision=3}]
  table[x=qlike,y=k]{figdata/methods/18-volatility-models/eval.csv};
\end{axis}
\end{tikzpicture}

\omcaption{Précision hors échantillon de trois prévisions de variance à un jour pour EUR/USD, 2015--2026 (3\,002 jours), en QLIKE moyenne relative à la fenêtre de 250 jours (une
valeur plus basse étant meilleure)~: GARCH-$t$ estimé sur 1999--2014 $-0.086$, EWMA $-0.060$. Les deux écarts sont significatifs selon le \omterm{def:qm:volatility-models:eval}{test de Diebold--Mariano}. Données~: taux de change de référence
de l'euro de la BCE (source~: statistiques de la BCE).}\label{fig:qm:volatility-models:eval}
\end{omfigure}

\section{Tutoriel~: le fantôme du rapport de risque}\label{tut:qm:volatility-models:ghost}

\begin{tutorial}
\textbf{Objectif.} Ajuster un GARCH sur EUR/USD, repérer le jour où un rendement important sort de la fenêtre de 250 jours, et comparer trois prévisions hors échantillon.
\textbf{État final~:} \cref{fig:qm:volatility-models:impulse,fig:qm:volatility-models:ghost,fig:qm:volatility-models:eval}, la \omterm{def:qm:stochastic-differential-equations:ou}{demi-vie} de 445 jours et la baisse
de 7.8~\%.
\begin{tutsteps}
\item \textbf{La récursion GARCH et sa vraisemblance}, avec des innovations gaussiennes ou de Student-$t$.
\omcode{code/firm/volfcst/firm_volfcst.py}{25}{38}{La récursion de variance GJR-GARCH et la log-vraisemblance quotidienne.}\label{lst:qm:volatility-models:filter}
\item \textbf{Fonctions de perte et \omterm{def:qm:volatility-models:eval}{test de Diebold--Mariano}.}
\omcode{code/firm/volfcst/firm_volfcst.py}{186}{215}{QLIKE, erreur quadratique, Mincer--Zarnowitz et Diebold--Mariano.}\label{lst:qm:volatility-models:dm}
\item \textbf{Exécuter} \texttt{fits()}, \texttt{ghost()}, \texttt{evaluate()} et \texttt{har\_vs\_ar1()} dans \texttt{qm\_vol.py}, puis \texttt{fig\_vol.py}.
\end{tutsteps}
\textbf{Pour aller plus loin.} Remplacer la fenêtre de 250 jours par une fenêtre qui pondère linéairement les rendements à mesure qu'on s'éloigne du bord, et observer la disparition du fantôme~; évaluer les prévisions
sur un horizon de dix jours, où le \omterm{def:qm:stochastic-differential-equations:ou}{retour à la moyenne} du GARCH devrait compter davantage.
\end{tutorial}

\section{Construction~: le module de prévision de volatilité}\label{bld:qm:volatility-models:volfcst}

\begin{build}
\bfield{Rôle} Toute volatilité que la firme miniature utilise pour le dimensionnement, les limites et les cotations d'options est une prévision issue de ce module, évaluée par rapport aux autres
chaque mois.
\bfield{Interface} \texttt{garch\_filter}, \texttt{garch\_fit(r, dist, gjr)}, \texttt{garch\_forecast(fit, horizon)}~; \texttt{ewma(r, lam)}, \texttt{rolling\_var(r,
window)}~; \texttt{har\_fit(rv)}, \texttt{har\_forecast(fit, rv)}~; \texttt{qlike}, \texttt{mse}, \texttt{mincer\_zarnowitz}, \texttt{diebold\_mariano}.
\bfield{Règles} Les prévisions du jour $t$ n'utilisent que les données jusqu'à $t - 1$~; les \omterm{def:qm:estimation:estimator}{erreurs types} sont celles du sandwich~; la persistance et la \omterm{def:qm:stochastic-differential-equations:ou}{demi-vie} sont indiquées avec chaque ajustement~; les comparaisons de prévisions
utilisent QLIKE ou l'erreur quadratique par rapport à un proxy sans biais, avec un \omterm{def:qm:volatility-models:eval}{test de Diebold--Mariano}.
\bfield{Tests d'acceptation} \texttt{code/firm/volfcst/tests/}~: GARCH et GARCH-$t$ retrouvent les paramètres simulés~; les récursions de GARCH, GJR, EWMA et de la fenêtre
calculées à la main~; le \omterm{def:qm:stochastic-differential-equations:ou}{retour à la moyenne} de la prévision de variance~; QLIKE est minimisée par la vraie variance sous un proxy bruité~; la pente de Mincer--Zarnowitz vaut un pour la vraie
variance~; le \omterm{def:qm:volatility-models:eval}{test de Diebold--Mariano} détecte une perte décalée~; HAR retrouve ses coefficients.
\bfield{Pour aller plus loin} EGARCH~; GARCH avec la \omterm{def:qm:volatility-models:rv}{variance réalisée} comme régresseur supplémentaire~; évaluation QLIKE sur plusieurs jours avec horizons qui se chevauchent.
\end{build}

\begin{omsources}
\item US Bureau of Labor Statistics, \emph{Consumer Price Index -- October 2022}, communiqué de presse du 10 novembre 2022~; BCE, taux de change de référence de l'euro (concertation vers 14~h~10 CET).
\item R.~F.~Engle, ``Autoregressive conditional heteroscedasticity with estimates of the variance of United Kingdom inflation'', \emph{Econometrica} 50, 1982.
\item T.~Bollerslev, ``Generalized autoregressive conditional heteroskedasticity'', \emph{Journal of Econometrics} 31, 1986.
\item L.~R.~Glosten, R.~Jagannathan et D.~E.~Runkle, ``On the relation between the expected value and the volatility of the nominal excess return on stocks'',
\emph{Journal of Finance} 48, 1993.
\item J.P.~Morgan/Reuters, \emph{RiskMetrics --- Technical Document}, 4th ed., 1996.
\item T.~G.~Andersen, T.~Bollerslev, F.~X.~Diebold et P.~Labys, ``Modeling and forecasting realized volatility'', \emph{Econometrica} 71, 2003~;
O.~E.~Barndorff-Nielsen et N.~Shephard, \emph{Journal of the Royal Statistical Society B} 64, 2002.
\item F.~Corsi, ``A simple approximate long-memory model of realized volatility'', \emph{Journal of Financial Econometrics} 7, 2009.
\item A.~J.~Patton, ``Volatility forecast comparison using imperfect volatility proxies'', \emph{Journal of Econometrics} 160, 2011.
\item F.~X.~Diebold et R.~S.~Mariano, ``Comparing predictive accuracy'', \emph{Journal of Business and Economic Statistics} 13, 1995~; J.~Mincer et
V.~Zarnowitz, ``The evaluation of economic forecasts'', NBER, 1969.
\item Banque centrale européenne, taux de change de référence de l'euro contre le dollar (ECB Data Portal, série \texttt{EXR.D.USD.EUR.SP00.A}), consulté le 24
septembre 2026.
\end{omsources}

\section{Exercices}

\begin{exercise}[$\star$]\label{exo:qm:volatility-models:1}
Un GARCH(1,1) a $\omega = 0.02$, $\alpha = 0.05$ et $\beta = 0.93$ (quotidien, rendements en pourcentage). Calculer ses volatilités inconditionnelles quotidienne et annuelle et
la \omterm{def:qm:stochastic-differential-equations:ou}{demi-vie} d'un choc de variance.
\end{exercise}

\begin{exercise}[$\star$]\label{exo:qm:volatility-models:2}
Avec $\lambda = 0.94$, au bout de combien de jours le poids d'un rendement au carré dans l'EWMA est-il divisé par deux~? Calculer la part du poids total portée par les 20 derniers jours.
\end{exercise}

\begin{exercise}[$\star$]\label{exo:qm:volatility-models:3}
\omterm{def:qm:volatility-models:rv}{Variance réalisée} à partir de rendements à 5 minutes sur 24 heures~: calculer son \omterm{def:qm:estimation:estimator}{erreur type} relative en tant qu'\omterm{def:qm:estimation:estimator}{estimateur} de la \omterm{def:qm:volatility-models:rv}{variance intégrée}, sans bruit.
\end{exercise}

\begin{exercise}[$\star\star$]\label{exo:qm:volatility-models:4}
Dans le même GARCH(1,1), la variance conditionnelle d'aujourd'hui vaut quatre fois la variance inconditionnelle. Calculer la prévision de la variance moyenne sur les 10 prochains jours.
\end{exercise}

\begin{exercise}[$\star\star$]\label{exo:qm:volatility-models:5}
Montrer que QLIKE, $\hat\sigma^2/h + \ln h$, est minimisée en espérance pour $h = \E[\hat\sigma^2]$, puis comparer les pertes des prévisions $0.5\sigma^2$ et $1.5\sigma^2$.
\end{exercise}

\begin{exercise}[$\star\star$]\label{exo:qm:volatility-models:6}
Une fenêtre de 250 jours contient un rendement de 3.5~\% parmi des jours à 0.5~\%. De quel pourcentage la volatilité de la fenêtre baisse-t-elle lorsque ce rendement en sort~?
\end{exercise}

\begin{exercise}[$\star\star\star$]\label{exo:qm:volatility-models:7}
\emph{Programmation.} Ajuster un GARCH(1,1) aux rendements EUR/USD avec la vraisemblance normale et indiquer les \omterm{def:qm:estimation:estimator}{erreurs types} de
$\alpha$ et de $\beta$, calculées à partir de la matrice hessienne puis de l'\omterm{def:qm:estimation:estimator}{estimateur} sandwich. Lesquelles faut-il rapporter, et pourquoi diffèrent-elles~?
\end{exercise}

\begin{exercise}[$\star\star\star$]\label{exo:qm:volatility-models:8}
\emph{Repérer l'erreur.} «~Nous avons comparé nos prévisions de volatilité par leur erreur absolue moyenne par rapport au rendement quotidien absolu~; l'EWMA avec $\lambda = 0.97$ l'emporte.~»
\end{exercise}

\section{Problème~: le fantôme du rapport de risque}

\begin{problem}[{Problème du week-end --- un choc de volatilité et le calendrier}]\label{pb:qm:volatility-models:1}
Les données sont les taux de change de référence quotidiens de l'euro contre le dollar publiés par la BCE, du 4 janvier 1999 au 23 septembre 2026, en log-rendements en pourcentage (7\,098 jours).

\medskip\noindent\textbf{Partie I --- L'ajustement.}
\begin{enumerate}
\item Quels paramètres GARCH(1,1) donne la vraisemblance normale, et quelles sont leurs \omterm{def:qm:estimation:estimator}{erreurs types} sandwich~?
\item Et avec la vraisemblance de Student-$t$~? De combien sa log-vraisemblance est-elle supérieure~?
\item Quelles sont la persistance et la \omterm{def:qm:stochastic-differential-equations:ou}{demi-vie} d'un choc de variance~?
\item Quelle est la volatilité quotidienne inconditionnelle, et comment se compare-t-elle à l'écart type empirique~?
\item Observe-t-on un \omterm{def:qm:volatility-models:gjr}{effet de levier}~?
\end{enumerate}

\medskip\noindent\textbf{Partie II --- Le fantôme.}
\begin{enumerate}[resume]
\item Quel rendement provoque la plus forte baisse relative sur un jour de la volatilité sur 250 jours, et à quelle date~?
\item De combien la volatilité de la fenêtre baisse-t-elle, et qu'a fait le marché ce jour-là~?
\item Que sont devenues les volatilités GARCH et EWMA le même jour~?
\item Combien de temps un rendement de 3.5~\% reste-t-il dans chaque prévision~?
\item Qu'aurait dit un rapport fondé sur la fenêtre au sujet du risque ce jour-là~?
\end{enumerate}

\medskip\noindent\textbf{Partie III --- Quelle prévision~?}
\begin{enumerate}[resume]
\item Quelles sont les \omterm{def:qm:volatility-models:eval}{pertes QLIKE} moyennes des trois prévisions de 2015 à 2026~?
\item Les écarts sont-ils significatifs~?
\item Que disent les pentes de Mincer--Zarnowitz~?
\item Pourquoi le $R^2$ de chaque prévision est-il si faible~?
\item Que changerait la \omterm{def:qm:volatility-models:rv}{variance réalisée} à partir de données intrajournalières dans l'évaluation~?
\end{enumerate}

\medskip\noindent\textbf{Partie IV --- Jugement.}
\begin{enumerate}[resume]
\item Quelle volatilité le rapport de risque doit-il utiliser~?
\item Une \omterm{def:qm:stochastic-differential-equations:ou}{demi-vie} de plus d'un an est-elle crédible~?
\item Que diriez-vous au comité des risques au sujet du jour du fantôme~?
\item Énoncer le \emph{résultat nommé}~: la \omterm{def:qm:stochastic-differential-equations:ou}{demi-vie} d'un choc de volatilité EUR/USD et l'ampleur du fantôme.
\item En une phrase~: qu'a de mauvais une fenêtre à pondération égale~?
\end{enumerate}
\end{problem}

\section{Questions d'entretien}

\begin{interviewq}[$\star$ \iqroles{risque, trader}]\label{iq:qm:volatility-models:1}
Pourquoi une volatilité historique à pondération égale peut-elle baisser un jour de forte volatilité~?
\end{interviewq}

\begin{interviewq}[$\star\star$ \iqroles{chercheur, risque}]\label{iq:qm:volatility-models:2}
Écrivez le GARCH(1,1). Quelle est sa variance inconditionnelle, et que signifie $\alpha + \beta$~?
\end{interviewq}

\begin{interviewq}[$\star\star$ \iqroles{chercheur}]\label{iq:qm:volatility-models:3}
Quel est le lien entre l'EWMA et le GARCH, et que fait-elle mal pour les prévisions à long horizon~?
\end{interviewq}

\begin{interviewq}[$\star\star$ \iqroles{chercheur, trader}]\label{iq:qm:volatility-models:4}
Comment évaluez-vous une prévision de volatilité lorsque la volatilité réelle n'est jamais observée~?
\end{interviewq}

\begin{interviewq}[$\star\star$ \iqroles{chercheur, trader}]\label{iq:qm:volatility-models:5}
Qu'est-ce que la \omterm{def:qm:volatility-models:rv}{variance réalisée}, et pourquoi ne pas l'échantillonner chaque seconde~?
\end{interviewq}

\begin{interviewq}[$\star\star\star$ \iqroles{chercheur}]\label{iq:qm:volatility-models:6}
Pourquoi le \omterm{def:qm:volatility-models:har}{modèle HAR} prévoit-il bien la volatilité avec seulement trois régresseurs~?
\end{interviewq}
